\section{本章小结}
\label{sec:metric-summary}

本章开头的问题是：在拓扑只告诉我们“哪些点在附近”的基础上，距离还能多说什么。第一个回答是把邻近变成可以计算的误差。度量的三条公理保证开球生成拓扑，这个拓扑满足Hausdorff条件，每点有可数邻域基。

邻域、闭包、收敛都可以用球和实数列的极限判断：闭包恰是到集合距离为零的点，也恰是集合中点列的极限。可数邻域基还使可分与第二可数在度量空间中等价，这一点在一般拓扑空间中并不成立。

第二个回答是区分距离带来的、拓扑看不出的性质。有界性和完备性都依赖所选度量：截断度量让整条实轴有界，开区间与实轴拓扑相同却不完备。完备性回答柯西列能否在空间内找到极限；在完备空间中，子空间完备当且仅当它是闭集；不完备的空间可以通过把柯西列识别为新点来完备化。

完备性还给出Baire纲定理：可数个稠密开集的交仍稠密，等价地，可数个无处稠密闭集不能覆盖空间，由此可以从逐点有界得到某个球上的统一有界。

第三个回答是把紧性拆成可检验的条件。度量空间紧致当且仅当序列紧致，也当且仅当完备且全有界：全有界保证总能选出柯西子列，完备性保证它有极限。有界闭集即紧集只是欧氏空间的特殊情形，无限离散空间是它在一般度量空间中的反例。子集相对紧当且仅当其中每个点列都有在环境中收敛的子列；环境完备时，这又等价于全有界。

最后，覆盖数、打包数和度量熵把全有界从“有限”细化为“多少”。球心取在环境中还是集合中、计数覆盖还是计数分离点，彼此只差一倍尺度；有限维赋范空间的半径$R$球需要约$(R/\varepsilon)^n$个半径$\varepsilon$的球，指数就是实维数。

这些结论都只涉及空间本身。第~\ref{chap:metric-space-maps}章转向空间之间的映射，研究误差经映射怎样传递，并把完备性用于稠密子集上的延拓和压缩迭代；第~\ref{chap:function-class-covers}章再把有限覆盖的计数用于函数类。
